% ======================================================================
% ΣΥΝΤΟΜΟ ΔΕΙΓΜΑ ΣΥΝΕΡΓΑΤΩΝ — Η ΙΔΙΑ ΥΛΗ ΣΕ SCREEN ΚΑΙ PRINT.
% Δοκιμαστική ύλη μόνο· δεν αποτελεί τελικό εγκεκριμένο κεφάλαιο.
% ======================================================================
\mainmatter
\BookChapterPhoto{Κίνηση σε μία διάσταση}{Δείγμα νέας εκδοτικής μακέτας — Φυσική Α΄ Λυκείου}
{assets/chapter-photo-sample.png}{Προσωρινό εικαστικό δείγμα για αξιολόγηση}{%
\item Να διακρίνουμε τη θέση, τη μετατόπιση και το διάστημα.
\item Να ορίζουμε τη θετική φορά του άξονα.
\item Να υπολογίζουμε τη μετατόπιση με μαθηματική σχέση.
\item Να ερμηνεύουμε την κλίση ενός διαγράμματος θέσης–χρόνου.
}
\BookSection{Θεωρία}
\section{Θέση και μετατόπιση}
Για να περιγράψουμε την κίνηση ενός σώματος ορίζουμε έναν άξονα \(x'x\) και επιλέγουμε τη θετική φορά. Η \textbf{μετατόπιση} ενός κινητού κατά τη μετακίνησή του από τη θέση \(x_A\) στη θέση \(x_B\) είναι:
\begin{KeyFormula}
\[ \Delta x=x_B-x_A \]
\end{KeyFormula}
\begin{center}
\begin{tikzpicture}[x=.85cm,>=Stealth]
  \draw[BookPrimary,->](-4,0)--(6,0) node[right] {\(x\)};
  \foreach \x in {-3,0,5}{\draw[BookPrimary](\x,0.08)--(\x,-0.08) node[below]{\scriptsize \(\x\)};}
  \fill[BookInk] (-3,0) circle(1.8pt) node[above]{\small A};
  \fill[BookInk] (5,0) circle(1.8pt) node[above]{\small B};
  \draw[BookSecondary,very thick,-{Stealth[length=3mm]}](-3,.7)--(5,.7) node[midway,above]{\(\Delta x\)};
\end{tikzpicture}
\end{center}
\begin{AttentionBox}
Η \textbf{μετατόπιση} εξαρτάται μόνο από την αρχική και την τελική θέση, ενώ το \textbf{διάστημα} εξαρτάται από τη διαδρομή.
\end{AttentionBox}
% Δείγματα ιστορικού ενθέτου και ρήσης. Χωρίς αποδιδόμενη πραγματική ρήση.
\begin{HistoryBox}[Η εξέλιξη της μέτρησης]
Η μέτρηση των φυσικών μεγεθών συνδέεται με την ιστορία των επιστημών.
Αυτό είναι δοκιμαστικό κείμενο που θα αντικατασταθεί από τεκμηριωμένη ιστορία.
\end{HistoryBox}
\BookQuote{Η παρατήρηση γεννά ερωτήσεις· η μέτρηση μας βοηθά να τις ελέγξουμε.}
{Δοκιμαστική ρήση}[Πρωτότυπο κείμενο επίδειξης μορφοποίησης]
% Κρατάμε τη λωρίδα μαζί με το πρώτο λυμένο παράδειγμα.
\Needspace{12\baselineskip}
\BookSection{Λυμένα παραδείγματα}
\begin{workedexample}{Μετατόπιση κινητού σε ευθεία}
\begin{ProblemBox}
Ένα κινητό βρίσκεται στη θέση \(x_A=-3\,\mathrm{m}\) και στη συνέχεια στη θέση \(x_B=5\,\mathrm{m}\). Να υπολογίσετε τη μετατόπιση και να προσδιορίσετε τη φορά της.
\end{ProblemBox}
\begin{ThinkBox}
Υπολογίζουμε τη διαφορά τελικής και αρχικής θέσης. Το πρόσημο του αποτελέσματος δηλώνει τη φορά ως προς τον άξονα.
\end{ThinkBox}
\SolutionHeading
\[\Delta x=x_B-x_A=5-(-3)=+8\,\mathrm{m}.\]
\end{workedexample}
\BookSection{Ερωτήσεις θεωρίας}
\BookQuestionHeading{Σωστό / Λάθος}
\begin{enumerate}[leftmargin=6mm]
\TrueFalseItem{Η μετατόπιση είναι διανυσματικό μέγεθος.}
\TrueFalseItem{Το διάστημα μπορεί να έχει αρνητική τιμή.}
\end{enumerate}
\BookSection{Άλυτες ασκήσεις}
\begin{exercise}[PHY-A-001]{}
Κινητό μετακινείται από τη θέση \(x_A=2\,\mathrm{m}\) στη θέση \(x_B=-7\,\mathrm{m}\). Να βρείτε τη μετατόπισή του και να εξετάσετε το πρόσημό της.
\end{exercise}
\begin{exercise}[PHY-A-002]{}
Ένα κινητό κινείται με σταθερή ταχύτητα \(\upsilon=4\,\mathrm{m/s}\) επί \(\Delta t=5\,\mathrm{s}\). Να υπολογιστεί η μετατόπισή του.
\end{exercise}
\begin{exercise}[PHY-A-003]{}
Κινητό μεταβαίνει από το σημείο \(A\) με \(x_A=0\,\mathrm{m}\) στο σημείο \(B\) με \(x_B=6\,\mathrm{m}\) και κατόπιν επιστρέφει στο \(C\) με \(x_C=2\,\mathrm{m}\).
\begin{enumerate}[label=\BookGreek*)]
\item Να υπολογίσετε τη συνολική μετατόπιση.
\item Να βρείτε το συνολικό διάστημα.
\end{enumerate}
\end{exercise}
\begin{exercise}[PHY-A-004]{}
Για κινητό που κινείται με σταθερή ταχύτητα δίνεται η σχέση \(x=2+3t\) (S.I.). Να προσδιορίσετε την αρχική θέση, την ταχύτητα και τη θέση τη χρονική στιγμή \(t=4\,\mathrm{s}\).
\end{exercise}
\clearpage
\BookSection{Λύσεις ασκήσεων}
\BookMinorHeading{Ενότητα 1 — Κίνηση σε μία διάσταση}
\begin{booksolution}{PHY-A-001}
Από τον ορισμό της μετατόπισης:
\[\Delta x=x_B-x_A=-7-2=-9\,\mathrm{m}.\]
Το αρνητικό πρόσημο δείχνει μετατόπιση προς την αρνητική φορά του άξονα.
\end{booksolution}
\begin{booksolution}{PHY-A-002}
Στην ευθύγραμμη ομαλή κίνηση:
\[\Delta x=\upsilon\Delta t=4\cdot5=20\,\mathrm{m}.\]
\end{booksolution}
\begin{booksolution}{PHY-A-003}
\textbf{α)} Η συνολική μετατόπιση εξαρτάται μόνο από την αρχική και την τελική θέση:
\[\Delta x=x_C-x_A=2-0=2\,\mathrm{m}.\]
\textbf{β)} Το συνολικό διάστημα είναι το άθροισμα των επιμέρους αποστάσεων:
\[s=|6-0|+|2-6|=6+4=10\,\mathrm{m}.\]
\end{booksolution}
\begin{booksolution}{PHY-A-004}
Συγκρίνοντας τη σχέση \(x=2+3t\) με τη \(x=x_0+\upsilon t\), βρίσκουμε:
\[x_0=2\,\mathrm{m},\qquad\upsilon=3\,\mathrm{m/s}.\]
Για \(t=4\,\mathrm{s}\): \(x=2+3\cdot 4=14\,\mathrm{m}\).
\end{booksolution}
